\section{Introduction}
A causal experiment supplies observations, admits specified controls and
carries an actual terminal mark. A reduction of that experiment consumes
its observations through a retained register. Its resource restriction
is not determined by the set of its output laws alone: the permitted
state at each cut and the visibility of persistent randomization must
also be specified. A finite-data auditor has its own such restrictions.
In particular, a test given an exact candidate law and a test learning
from repeated trials solve different mathematical problems.

The object organizing this paper is a family of continuation responses.
At a causal cut, those responses must be compatible with the one-step
shifts at the next cut. At a training--validation cut, they are the
conditional validation responses available from each retained state.
The first viewpoint characterizes exact realization; the second gives a
geometric spectrum for approximate statistical certification. We show
that these viewpoints control the same physical audit: a lower bound on
its possible continuation responses, an attaining finite-sample rule,
and the exact continuation profile of that rule.

\subsection{The decision cut and the principal theorem}
Let $B$ be a compact set of laws on a finite marked alphabet $\Omega$,
and let $Q$ be the target law. A reset auditor trains on independent
samples from one candidate $b\in B$, held fixed for the whole audit.
Only its charged register survives the end of training. It then selects
an event before independent target and candidate validations and emits
the difference of their event indicators. Its score means the expectation
of this signed difference. A response $h\in[0,1]^\Omega$ is the mean
indicator of a randomized event. Put
\[
 \mathfrak V_K(B,Q)=\max_{h_1,\ldots,h_K}
          \min_{b\in B}\max_{m\le K}(Q-b)h_m.
\]
This is a continuation invariant of the marked decision interface.
It does not convexify the candidate image and does not grant the auditor
an external copy of its training record.

\begin{theorem}[Decision continuation complexity and physical certification]
\label{thm:v20-main}
The following assertions hold with all persistent data charged to the
appropriate machine.

\smallskip\noindent
\textup{(i)} The supremum $v_{n,K}$ of worst-case expected scores with
$n$ training samples and $K$ states at the decision cut satisfies
\[
 0\le\mathfrak V_K-v_{n,K}
 \le\min\left\{\sqrt{(|\Omega|-1)/n},
                    \sqrt{2\log(2K)/n}\right\}.
\]
The least $K$ permitting a strict score greater than $c$ at some finite
training length is the least size of a cover of $B$ by sets
$\{b:(Q-b)h>c\}$. This spectrum is Lipschitz under uniform total
variation perturbations of candidate and target laws.

\smallskip\noindent
\textup{(ii)} For the ideal marked private product family $B$ of the
two-sphere experiment, $\mathfrak V_1=1/8$ and
$\mathfrak V_2=3/8$. Throughout the collision interval
$9/10\le d\le11/10$, $0<\sigma\le1/24$, the minimum decision-cut
width for uniform physical score greater than $2/5$ is exactly three.
The lower bound applies to every finite-reset auditor with the specified
interface, not merely a particular counter implementation; feedback-row selection
does not evade it.

\smallskip\noindent
\textup{(iii)} One deterministic rational rule has physical score greater
than
\[
 \sqrt2-1-\frac1{4\sqrt{399e}}-\frac14e^{-399/8}-\frac1{300}
                         >\frac25.
\]
It uses $399$ informative training trials, an optional ignored training
trial and two independent validations, hence fits the $402$-trial
schedule. It has exactly three decision-cut states. On its informative
training segment, the minimum simultaneous widths for exact stochastic
implementation are
\[
 r_t=\begin{cases}
 \binom{t+2}{2},&0\le t\le199,\\
 \binom{402-t}{2},&200\le t\le399.
 \end{cases}
\]
Their maximum is $20301$. Including individual report-bit acquisition
and the validations, at most $40602<2^{16}$ retained states suffice.

\smallskip\noindent
\textup{(iv)} At an arbitrary cut of a finite-valued decision, a weighted
graph of distinguishing suffixes gives a distribution-sensitive error
bound through its minimum monochromatic weight under $K$ labels.
Applied to Bayes margins, it yields a quantitative constrained-audit
regret bound. At zero error it recovers the exact simultaneous residual
profile, even for stochastic implementations.

\smallskip\noindent
\textup{(v)} For fixed positive $\sigma$ and changing $d$, the same
physical rule has the target-variation modulus
$\min\{1,\sqrt{|d'-d|+2700|d'-d|^2}/(2\sigma)\}$, with no training-count
factor. The microscopic paths, likelihoods, posterior processes,
innovations, and bounded entropic backward values, integrands and
martingale integrals have the joint transport conclusions proved below.
The private and visible deficiencies exceed $9/20$, whereas a hidden
two-valued selector has error less than $13/100$. In the no-collision
regime $d\le2/5$ all three finite-experiment deficiencies vanish.
\end{theorem}
\begin{proof}
Part (i) is Theorem~\ref{thm:v20-spectrum}. The exact ideal spectral
values are Theorem~\ref{thm:v20-two}; Lemmas~\ref{lem:v20-envelope}
and~\ref{lem:v20-majority} and Theorem~\ref{thm:v20-physical} prove
the matching physical construction. Theorem~\ref{thm:v20-profile}
proves (iii), including all acquisition and validation phases.
Theorem~\ref{thm:v20-weighted} and Corollary~\ref{cor:v20-bayes}
prove (iv). The microscopic estimates in
Theorem~\ref{thm:v18-microscopic}, the explicit coupling-version lemmas
in Section~\ref{sec:v20-versions}, and
Theorem~\ref{thm:v19-global} prove (v).
\end{proof}

The key point is that the two-state obstruction and the three-state
construction are for the same physical score. Three witness candidates
force three validation continuations above a prescribed level; a
three-event dictionary then covers the entire original nonconvex product
image. Its empirical majority selector has a residual quotient that can
be computed exactly at every training cut. Thus the continuation invariant
both limits the physical certificate and constructs it. The larger
unrestricted marked discrepancy is a different optimization. Likewise,
minimum width at one specified cut and minimum peak width over all
successful algorithms are different quantities. The theorem gives an
exact answer to the former, and an exact implementation profile for the
stated successful rule.

\subsection{The full comparison and transport framework}
The general finite experiment has a private register profile $K$, a
possible independent hidden selector, or a visible selector included in
the compared law. Compatible normalized continuation covers characterize
its exact causal realization. Compactness, nonnegative factorization,
priced composition and hidden/visible comparison are proved in
Theorems~\ref{thm:v13-main} and~\ref{thm:v18-continuation}. The actual
mark cannot be read by the online candidate. The reset auditor's completed
marked trials are a separate interface with a separate paid register.

The finite-reset minimax theorem and its moment representation operate on
the original compact behavior image, without replacing a sampled
candidate by the barycenter of a prior. The sampling error is uniformly
of order $n^{-1/2}$, with a matching binary sampling obstruction. The
Boolean streaming theorem determines the widths needed to attain that
binary testing value. The weighted finite-valued result extends its
error mechanism to nonuniform word weights and general finite decisions.
The nonlinear revealed-behavior witness, Bernstein hierarchy and rational
chart retain their full proofs. Their evaluator sees an exact behavior
law, unlike a reset auditor; their privileges are not interchanged.

The physical model consists of two labeled hard spheres with a prepared
impact--mark correlation and one nongrazing collision. The collision
transduces the impact variable into an observed transverse velocity;
the acquisition gate does not change the mechanics. Varying the contact
diameter changes the collision time and trajectory. A common-preparation
Gaussian observation estimate transports both the finite certificate and
its conditional-law consumer. The latter is a bounded scalar entropic
backward equation with fixed positive observation noise. Its terminal
condition is the observation-measurable conditional certainty equivalent,
not an unobserved latent payoff. The bounded tilt yields global
integrand transport without a density stopping time. The coupled
stochastic integrals retain their own marginal observation filtrations.

\subsection{Relation to classical comparison and realization}
The classical comparison theory of Blackwell and Le Cam
\cite{Blackwell1953,LeCam1964,Torgersen1991} and earlier filtered
comparison \cite{Norberg2002,Weisshaupt2006} provide essential context.
Introducing a filtration or a randomization criterion is not itself a
novelty claim. Stable positive-cone realization \cite{Fliess1975} and
residual-state minimization \cite{Nerode1958} are classical mechanisms.
Their finite controlled normalization is proved here so that the resource
profile and the input interface remain explicit. We do not claim a new
deterministic automata minimization theorem.

The contribution is the link from compatible continuation responses to
statistical certification at a charged physical decision cut. The
response-envelope classification has a uniform finite-sample realization;
its first two values on the physical product image can be computed
exactly; and a three-response cover attains the prescribed microscopic
score with a completely specified streaming quotient. The weighted
obstruction uses state-distribution overlap and concavity rather than an
unsupported deterministic reduction of randomized memory. Scattering,
Bayes' formula, stochastic exponentials and martingale inequalities are
used as tools, not claimed as discoveries.

The original full text of Norberg's paper was not obtained for a
proof-level comparison. The literature crosswalk therefore separates
verified bibliographic information, inherited original-source comparisons,
and unresolved priority questions. In particular, no absence claim about
Norberg's theorems is inferred from this access limitation. The results
here are presented by their precise statements and proofs, not an
unverified assertion of priority for filtered comparison.

The larger theta-theory program retains all eleven historical components.
The program map in Appendix~\ref{sec:v20-program} distinguishes theorem
inputs, scoped replacements, historical aggregate obligations and the
independent A2 geometric chain. The present foundational layer organizes
the causal comparison and certification chain; independence of another
chain is not replaced by a nominal dependency arrow. All predecessor
proofs and program targets remain available in the unchanged historical
companion to this article.
